\documentclass[11pt,a4paper]{article}
\usepackage[top=1.0in, bottom=1.0in, left=1.1in, right=1.1in]{geometry}
\usepackage{graphicx}
\usepackage[sort&compress, super, numbers]{natbib}
\usepackage[export]{adjustbox}
\usepackage[ngerman]{babel}
\usepackage[utf8]{inputenc}
\usepackage[T1]{fontenc}
\usepackage{parskip}
\setlength{\parindent}{0pt}
\renewcommand{\bibname}{References}%names reference list 

\begin{document}

\pagenumbering{gobble}
\includegraphics[width=0.4\textwidth, right]{ubcLetterHead.png}
Dear Dr. whatever
\vspace{1.5ex}\\
Please consider our paper, ``Warmer and longer growing seasons increase growth for juvenile trees in a common garden'' for publication as a full paper in \emph{New Phytologist}. 
\vspace{1.5ex}\\ 

Anthropogenic climate change affects the timing of many biological processes, producing cascading effects. For temperate trees, warming temperatures lengthen their growing season, which should increase growth. However, recent evidence showed that trees may not grow more with warmer seasons which could be detrimental for the critical role that temperate forest carbon sink has on regulating greenhouse gas emissions. Our `Full paper' adds a rare combination of full growing season length via vegetative phenology and growth via tree rings to address. The favourable context of our common garden pr this fundamental question  evidence that under drought- and competition-limited conditions, juvenile trees grew more with longer and warmer seasons. This adds important information to the important question that will determine whether forests regenerate and store carbon with future climate.
% Climate change is altering the environmental cues of growth globally, producing shifts in the timing of life history events (phenology), altering resource use and producing novel species interactions. 
% But to date, we have observed considerable variation in communities' responses to environmental change, likely reflecting differences in species assemblages and the effects of other species' traits on the timing and patterns of growth. In plants, decades of research has found consistent relationships between traits and growth strategies under different environments---giving rise to frameworks like the leaf or wood economic spectra \citep{Wright2004, Chave2009} that separate species along a gradient of fast and acquisitive to slow and conservative  strategies. Phenology was foundational to the theory that underlies these spectra \citep{Grime1977}, but has rarely been included because it is often challenging to measure and highly variable across species when measured in natural conditions. 
%  likely to correlate with leaf and plant traits associated with species timing and growth strategies.   
% 50 word answers

\textit{What hypotheses or questions does this work address?}  
Why trees do not grow more with longer and warmer seasons? This outstanding question, that we address here in a common garden using juvenile trees of four species sourced from four provenances. We test effects of four different growing season metrics on growth and determined whether provenance latitude changed growth. 

\textit{How does this work advance our current understanding of plant science?}  
Most recent evidence that growth may not increase with longer seasons was done in natural setting with many confounding factors, which made hard understanding the effect that growing season length had on growth. By a using multi-year, multi-species high quality phenology and growth dataset, we are able to tease apart the contribution of each season metric. We show in particular that the end of season events---most often overlooked---appeared to be more important for growth than spring events. 
% Don't repeat the abstract too much: By measuring 1428 individuals from forests spanning 6$^{\circ}$ latitude and 55$^{\circ}$ longitude ...
% We leverage massive experiments with new modeling approaches to show that phenology fits within major frameworks for plant growth, with early to late budburst correlating with plant size and leaf metrics. Further, we found photoperiod appears most linked to other plant traits, despite temperature being the dominant controller of leafout. 
\vspace{1.5ex}\\ 
%  Surprisingly, we found photoperiod to be the main cue linking phenology with broader plant growth strategies, suggesting trade-offs between traits and phenology may limit future climate responses.
\vspace{1.5ex}\\
\textit{Why is this work important and timely?}
Our work provides at a critical time in climate change ecology challenging recent findings on tree growth. We show that young deciduous trees may actually thrive under a warmer climate with  has critical implication for forest regeneration and range expansion demonstrating that growing season length and growth may not relate much. Further, we show that the end of season may be more important than the start of season.
%emw2026Apr12 -- Did you find support for early vs late? Do you overcome perhaps a major previous limitation in trait frameworks and show that phenology firmly belongs in the frameworks and is far less variable than thought? Remember, it's the cover letter -- go big!
% Our results provide the most compelling evidence of trait-phenology relationships across North America with implications for how forest communities will shift with climate change. Further, our novel approach shows how controlled environment experiments combined with cutting-edge hierarchical modeling can elucidate the role of phenology in structuring plant growth strategies. 
\vspace{1.5ex}\\ 
% We found that species phenologies partially do align with existing trait frameworks and differences in growth strategies, while also highlighting the importance of additional environmental cues or physiological trade-offs. Through our novel multi-trait approach, we can better predict species’ vulnerability to climate change and more accurately forecast future community responses.
\vspace{1.5ex}\\
All authors contributed to this work and approved this version for submission. The manuscript is 4083 words with a 200 word summary, and 4 figures and is not under consideration elsewhere. We hope you find it suitable for publication in \emph{New Phytologist}, and look forward to hearing from you.
Sincerely, \\
\includegraphics[scale=.2]{signature.pdf} \\
Christophe Rouleau-Desrochers\\
MSc in Forestry \\ 
Department of Forest \& Conservation Sciences \\
The University of British Columbia
\newpage
\vspace{-5ex}

\begingroup
\renewcommand{\section}[2]{}
\bibliographystyle{gcb.bst}% 
\bibliography{ExportedItems.bib}
\endgroup
\newpage

\end{document}


%%%%%%%%%%%%%%%%%%%%%%%%%%%%%%%%%%%%%%%%%%%%%%%%%%%%%%%%%%%%%%%%%%%%%%%%%%%%%%%%

%%%%%%%%%%%%%%%%%%%%%%%%%%%%%%%%%%%%%%%%%%%%%%%%%%%%%%%%%%%%%%%%%%%%%%%%%%%%%%%%

\textit{What scientific question is addressed in this manuscript?}  
Do trees grow more with longer, warmer seasons? To answer this question, we used a unique common garden dataset of radial growth data from juvenile trees of four species sourced from multiple provenances. Using paired leaf phenology data on budburst and budset we tested the effect of both longer and warmer seasons. 

\textit{What is/are the key finding(s) that answer this question?}
We found that warmer seasons increased tree growth, with thermal season being a better predictor for growth than calendar season. We suggest later season warmth may be more important to growth, which is supported by later budset driving increased growth and a little role of an earlier leafout.

\textit{Why is this work important and timely?}
Our work brings important insights on how longer, warmer seasons affect tree growth. Recent work has challenged the assumption that longer growing seasons drive increased tree growth, with major implications for the future of temperate forest carbon sink in a warming climate \citep{dow_warm_2022,silvestro_longer_2023,matula_shifts_2023}.  

\textit{Describe how your paper fits within the scope of GCB.}
Our work provides new evidence about how trees grow in response to longer and warmer seasons, with implications for carbon storage and forest dynamics. We show that young deciduous trees may actually thrive under a warmer climate, which could affect forest regeneration and range expansion with warming \citep{luu_importance_2024}.

\textit{What biological AND global change aspects does it address?}
Our paper addresses fundamental questions about how tree grow, from the angle of how climate change will alter growing seasons and thus tree growth. Our results suggest the role of late-season events and the thermal season may be critical to understanding and predicting tree growth in the future. 

\textit{What are the three most recently published papers that are relevant to this manuscript?}\\
\cite{dow_warm_2022}, \citep{silvestro_longer_2023}, \citep{wolkovich_why_2025}

\textit{Provide the name and institutional email addresses of five potential reviewers.} \\
Caroline Polgar---\href{mailto:carolinepolgar@gmail.com}{carolinepolgar@gmail.com} \\
Amanda Gallinat---\href{mailto:agallina@colby.edu}{agallina@colby.edu} \\
Libby Ellwood---\href{mailto:eellwood@floridamuseum.ufl.edu}{eellwood@floridamuseum.ufl.edu} \\
Loïc D'Orangeville---\href{mailto:loic.dorangeville@sbf.ulaval.ca}{loic.dorangeville@sbf.ulaval.ca} \\
Julia Green---\href{juliakgreen@arizona.edu}{juliakgreen@arizona.edu}

All authors contributed to this work and approved this version for submission. The manuscript is 3977 words with a 2820-word abstract, and 3 figures and is not under consideration elsewhere. We hope you find it suitable for publication in \emph{Global Change Biology}, and look forward to hearing from you.

